\documentclass[../../main.tex]{subfiles}
\begin{document}

\begin{flushright}
\footnotesize\textit{【自動文字起こし・要確認】}
\end{flushright}

{\bf [解]} $A, B$ の $n$ 回目の試行後の座標を $a_n, b_n$ とする．又以下のように場合の数を置く．

\begin{itemize}
  \item[$1^\circ$] $a_n > b_n \quad (P_n)$ \\
  表なら共に $+1$，裏なら $B$ を $+1$ すすめる．
  \item[$2^\circ$] $a_n = b_n \quad (X_n)$ \\
  表なら $A$ を $+1$ すすめる，裏なら $B$ を $+1$ すすめる．
  \item[$3^\circ$] $a_n < b_n \quad (Q_n)$ \\
  表なら $A$ を $+1$ すすめる，裏なら共に $+1$ すすめる．
\end{itemize}

\medskip

\noindent
\paragraph{(1)}

題意から，任意の $n$ に対して，$|a_n - b_n| \le 1$ である．したがって $n+1$ 回目に $a_{n+1} = b_{n+1}$ となるのは以下の場合：
\begin{align}
\begin{cases}
1^\circ \text{ で裏が出る．} \\
3^\circ \text{ で表が出る．}
\end{cases}
\end{align}
したがって，
\begin{align}
X_{n+1} &= P_n + Q_n = 2^n - X_n \quad (\because P_n + Q_n + X_n = 2^n) \quad \cdots \textcircled{1}
\end{align}

\medskip

\noindent
\paragraph{(2)}

$Y_n = \dfrac{X_n}{2^n}$ によって $\{Y_n\}$ を定める．\textcircled{1}の両辺を $2^{n+1}$ でわって
\begin{align}
Y_{n+1} &= -\dfrac{1}{2}Y_n + \dfrac{1}{2}
\end{align}
\begin{align}
Y_{n+1} - \dfrac{1}{3} &= -\dfrac{1}{2}\left(Y_n - \dfrac{1}{3}\right) \quad \cdots \textcircled{2}
\end{align}
$X_1 = 0$ つまり $Y_1 = 0$ だから，\textcircled{2}をくり返し用いて，
\begin{align}
Y_n - \dfrac{1}{3} &= \left(-\dfrac{1}{2}\right)^{n-1}\left(-\dfrac{1}{3}\right)
\end{align}
\begin{align}
\therefore X_n &= \dfrac{1}{3} \{ 2^n + 2 \cdot (-1)^n \} \hfill \text{\slash\slash}
\end{align}

\medskip

\noindent
\paragraph{(3)}

$A_n$ の期待値 $E_n(a)$ を求めれば良い．ここで，確率変数 $t_k$ を以下のように定める．
\begin{align}
t_k = \begin{cases}
0 & (k \text{回目の試行で } A \text{ が進まない}) \\
1 & (k \text{回目の試行で } A \text{ が進む})
\end{cases}
\end{align}
すると，題意から
\begin{align}
E_n(a) &= \sum_{k=1}^n P(t_k = 1) \quad \cdots \textcircled{3}
\end{align}
である．$k$ 回目で $A$ が進むのは以下の時：
\begin{align}
\begin{cases}
\text{\textcircled{1} コインで表が出る． (確率 } \dfrac{1}{2}\text{)} \\
\text{\textcircled{2} } a_{k-1} < b_{k-1} \text{ かつ コインで裏が出る．}
\end{cases}
\end{align}
ここで，(1)及び(2)，対称性から $P_n = Q_n$ であることから，
\begin{align}
P_n &= \dfrac{1}{2} (2^n - X_n) = \dfrac{1}{3} (2^n - (-1)^n)
\end{align}
だから，\textcircled{2}の起きる確率は $\dfrac{1}{2^k} P_{k-1}$ である．したがって，$k \ge 2$ の時，
\begin{align}
P(t_k = 1) &= \dfrac{1}{2} + \dfrac{1}{2^k} \cdot \dfrac{1}{3} ( 2^{k-1} - (-1)^{k-1} ) \\
&= \dfrac{2}{3} + \dfrac{1}{3} \left(-\dfrac{1}{2}\right)^k
\end{align}
である．$k=1$ の時，$P(t_1=1) = \dfrac{1}{2}$ であるからこれも良い．したがって\textcircled{3}から，
\begin{align}
E_n(a) &= \sum_{k=1}^n \left\{ \dfrac{2}{3} + \dfrac{1}{3}\left(-\dfrac{1}{2}\right)^k \right\} \\
&= \dfrac{2}{3}n - \dfrac{1}{6} \cdot \dfrac{1 - (-1/2)^n}{1 + 1/2} \\
&= \dfrac{2}{3}n - \dfrac{1}{9}\left\{ 1 - \left(-\dfrac{1}{2}\right)^n \right\} \hfill \text{\slash\slash}
\end{align}

\end{document}
